\begin{figure}[ht]
\centering
\begin{tikzpicture}[x=1.7cm,y=1.25cm,>=Stealth,
 every node/.style={font=\small}]
 \draw[gray!60] (0,0) rectangle (3,2.2);
 \draw[densely dashed,gray!75] (0,0)--(3,2.2)
 node[pos=.48,above left,text=black] {\(\Delta\)};
 \draw[very thick,blue!65!black,->] (0,2.2)--(2.8,2.2);
 \draw[very thick,blue!65!black,->] (3,0)--(3,2.05);
 \fill (0,0) circle (1.5pt);
 \fill (0,2.2) circle (1.5pt);
 \fill (3,0) circle (1.5pt);
 \fill (3,2.2) circle (1.8pt);
 \node[below left] at (0,0) {\((a,a)\)};
 \node[left] at (0,2.2) {\((a,b)\)};
 \node[below] at (3,0) {\((b,a)\)};
 \node[above right] at (3,2.2) {\((b,b)\)};
 \node[above] at (1.4,2.2) {\(L_1:\ s_2=b\)};
 \node[right] at (3,1.05) {\(L_2:\ s_1=b\)};
 \node at (1.15,.45) {\(F|_\Delta=g\)};
\end{tikzpicture}
\caption{The two factor curves used to test the bracket tensors,
drawn schematically in \(C\times C\).
On \(L_i\simeq\PP^1\), the input bundle is \(g^*E_i\) and the
complementary target bundle is trivial. Sections that vanish at
the marked off-diagonal point force their tensor values to
vanish at \((b,b)\), which maps to \(g(b)\).}
\label{fig:box}
\end{figure}

